%% kapa: exam=1 round=23 year=2012 period=1 form=A subject=civil_law subject_ko=민법 seq=9 number=9
%% source: Q-Net 감정평가사 기출문제 (cst00309) artlSeq=5052578 file=raw/1008/5052578/2012년 제23회 감정평가사 제1차시험 기출문제(1교시 A형).pdf
%% conversion: kordoc 4.24.51 -> tools/gen1.py (auto); review=proofread-sonnet; flags=none
\begin{question}{9}
X 부동산을 甲으로부터 매수한 乙은 이전등기를 경료하지 않은 채 丙에게 그 부동산을 매도하였고 현재 丙은 이를 점유하여 사용․수익하고 있다. 다음 설명 중 옳지 않은 것은? (다툼이 있으면 판례에 의함)
\begin{choices}
  \item 3자간 중간생략등기의 합의가 있더라도 甲과 丙 사이에 토지거래허가를 받아 甲에서 丙 앞으로 직접 경료된 소유권이전등기는 무효이다.
  \item 甲과 丙 사이에 물권적 합의가 없더라도 중간생략등기에 관한 3자 합의가 있다면, 丙은 甲에 대하여 직접 자신의 명의로 이전등기를 청구할 수 있다.
  \item 중간생략등기의 합의가 있은 후 甲과 乙 사이에 매매대금을 인상하는 약정이 체결된 경우, 인상된 매매대금이 지급되지 않았음을 이유로 甲은 丙에 대한 소유권이전등기의무의 이행을 거절할 수 없다.
  \item 이미 甲에서 丙으로 소유권이전등기가 경료된 경우, 중간생략등기에 관한 합의가 없었다는 이유만으로는 중간생략등기가 무효라고 할 수 없다.
  \item 乙의 甲에 대한 이전등기청구권은 중간생략등기의 합의가 있는 경우에도 소멸되지 않는다.
\end{choices}
\end{question}
\answer{3}
